Large language models (LLMs) trained with reinforcement learning from human feedback (RLHF) routinely produce gratitude-expressing text, yet it remains unclear whether this behavior reflects genuine preference-like states or sophisticated pattern matching on polite human data. We present three experiments probing LLM ``thankfulness'' along complementary dimensions: expression consistency, reception sensitivity, and self-report--behavior alignment. Using \gptfour, we find that (1) the model expresses gratitude with 95\% context-sensitivity, correctly distinguishing situations where gratitude is socially appropriate from those where it is not ($p < 0.001$, $\varphi = 0.804$); (2) user gratitude increases response length by 2.7$\times$ but does not improve---and may slightly decrease---response quality; and (3) self-reported gratitude shifts readily under persona manipulation while behavioral cooperation does not, yielding a non-significant self-report--behavior correlation ($r = 0.50$, $p = 0.67$). These findings suggest that LLM thankfulness is best characterized as a linguistically competent but behaviorally ungrounded performance: the model has learned \emph{when} to say ``thank you'' without developing corresponding preference states that influence downstream action. We discuss implications for AI interaction design, alignment evaluation, and the philosophy of artificial preferences.